\begin{houseviewbox}[重估不是“业绩披露后看股价”，而是逐项替换模型假设]
本报告当前的 17.0 元目标和减持/避免新增仓位，表达的是在 17.75 元现价下缺乏足够安全边际。正式 H1 披露后，应在同一张模型里更新已验证的 H1 数字、再重建 H2，而不是将预告兑现简单视为利好或利空。
\end{houseviewbox}

\section{重估顺序：先现金，后单位经济，再估值倍数}

最先需要确认的是利润到现金的转换。若 H1 归母处于预告区间，却仍由应收和存货扩张驱动，估值折价应保留；若经营现金改善并有营运资本解释，才可降低熊情景概率。第二步是验证锂盐销量、实现收入代理、毛利或可替代运营信息；第三步才是讨论资源供给、民爆项目和估值倍数。这个顺序避免在缺少基础经营数据时，先讨论题材或目标价。

\begin{exhibitbox}[表 13-1：正式 H1 的重估决策树]
\begin{tabularx}{\exhibitboxwidth}{L{2.8cm}X X X}
\toprule
检查项 & 上修信号 & 降级信号 & 模型动作 \\
\midrule
归母利润 & 落在预告中上部且经营桥可解释 & 低于预告或非经常项主导 & 更新 H1 锚，不年化 Q2 \\
经营现金流 & 转正或明显改善，应收/存货有闭环 & 继续明显为负且占款扩大 & 调整 OCF、FCF 和估值折价 \\
锂盐单位经济 & 销量、收入/吨、毛利支持 H2 基准 & 与 9 万吨、10.5 万元/吨、32\% 组合不符 & 下调 H2 锂收入和毛利 \\
民爆/矿服 & 利润与回款同步、分部信息完整 & 增长伴随项目占款或成本压力 & 维持或下调防御性利润池 \\
资源和项目 & 实际供应、成本或期限可复算 & 仍只有设计与意向信息 & 仅影响可选上行，不改基准 \\
\bottomrule
\end{tabularx}
\end{exhibitbox}

\section{三种披露结果与投资行动}

\begin{exhibitbox}[表 13-2：结果矩阵——从证据到观点]
\begin{tabularx}{\exhibitboxwidth}{L{2.2cm}X X X}
\toprule
结果 & 需同时满足的条件 & 概率/估值处理 & 行动 \\
\midrule
强验证 & 利润、现金、单位经济三者闭合；H2 有可观察支持 & 牛概率上调，复核 PE 与 PB/ROE 参数 & 由防守转中性，再评估新增仓位 \\
部分验证 & 利润达标，但现金或单位经济仍缺一环 & 保持基准或小幅调整，倍数不扩张 & 继续等待，不因标题性利好追价 \\
负验证 & 利润、现金或 H2 假设任一关键项显著失配 & 熊概率上调，重算 EPS 与价值区间 & 减持/退出至观察名单 \\
\bottomrule
\end{tabularx}
\end{exhibitbox}

\section{报告的可证伪性}

本观点不依赖“锂价继续上涨”这一难以审计的单句判断。它可以被三个方向证伪：第一，若正式 H1 表明现金流和营运资本显著好于基准，当前 15\% 市场锚与 85\% 内在价值的权重应上调基本面价值；第二，若实际单位经济使 H2 归母稳定在 11 亿元以上，牛情景应提高；第三，若现金继续恶化或 H2 无法支持 4 亿元归母，熊情景的 10.3 元价值应成为重新检验的起点。对机构投资者而言，能被新披露改变的结论，才是可管理的结论。

\sourcenote{House 风险框架、估值模型与 2026H1 业绩预告。}
